% Abhakseite „Das kann ich“ – Zone nur im Gesamt (\mitzone)
%% TODO Ich-kann-Sätze: Platzhalter wie in den Aufgabendateien
\begin{abhakseite}
\ifdefined\mitzone
\abhakgruppe{Kennst du schon}
\abhak{1}{Ergebnismenge eines ein- oder zweistufigen Versuchs als geordnete Paare, Laplace-Regel einstufig, Summenregel, Gegenereignis, Zufallsgerät zu einer Wahrscheinlichkeit entwerfen}
\abhak{2}{Möglichkeiten vollständig aufzählen, Zählprinzip „mal“, Anordnungen einer kleinen Menge}
\abhak{3}{Brüche multiplizieren (auch mehrere Faktoren mit fallenden Nennern), kürzen, gleichnamig und ungleichnamig addieren; Bruch ↔ Dezimalzahl ↔ Prozent; Potenzen kleiner Brüche}
\abhak{4}{Symbolschreibweise der Ereignisse (A ∩ B, A ∪ B, Ā, A \ B, $P_A(B))$ lesen; die bedingte Wahrscheinlichkeit als Quotient}
\abhak{5}{Anteil mal Gesamtzahl (absolute Anzahl), Mittelpunktswinkel als Anteil des Vollwinkels, Prozentsatz als Dezimalzahl}
\abhak{6}{Baumdiagramm mit Zurücklegen oder mit zwei Geräten zeichnen und ergänzen (Äste an einem Punkt ergeben eins), Pfadregel, Summenregel, „mindestens einmal“ über das Gegenereignis, dreistufig}
\abhak{7}{Ziehen ohne Zurücklegen zweistufig (Nenner um eins kleiner, bei der gezogenen Sorte auch der Zähler), zweite Person zieht, dreistufig ein Pfad, „unter den ersten drei“}
\abhak{8}{Kombinatorik}
\abhak{9}{Lineare Gleichungen, quadratische Gleichungen (Lösungsformel, Lösungen sieben) und einfache Bruchgleichungen aufstellen und lösen}
\abhak{10}{Ziehen ohne Zurücklegen zweistufig (Nenner um eins kleiner, bei der gezogenen Sorte auch der Zähler), zweite Person zieht, dreistufig ein Pfad, „unter den ersten drei“ – Fehler finden}
\abhak{11}{Ziehen ohne Zurücklegen zweistufig (Nenner um eins kleiner, bei der gezogenen Sorte auch der Zähler), zweite Person zieht, dreistufig ein Pfad, „unter den ersten drei“ – selbst rechnen}
\fi
\abhakgruppe{Einheit 1 · Ereignisse als Mengen}
\abhak{12}{Ereignisse als Mengen}
\abhak{13}{Ereignisse als Mengen – Fehler finden, Begründen, Anwendung, Darstellungswechsel}
\abhakgruppe{Einheit 2 · Laplace-Experimente in der Oberstufe}
\abhak{14}{Laplace-Experimente}
\abhak{15}{Laplace-Experiment: Aussage über Wahrscheinlichkeiten aus absoluten Anzahlen ohne Grundgesamtheit beurteilen}
\abhak{16}{Laplace-Experimente – Fehler finden, Begründen, Anwendung, Darstellungswechsel}
\abhakgruppe{Einheit 3 · Pfadregeln bei unabhängigen Stufen}
\abhak{17}{Pfadregeln bei unabhängigen Stufen}
\abhak{18}{Gleich wahrscheinliche Ergebnisse eines zweistufigen Experiments über die Pfadwahrscheinlichkeiten ermitteln}
\abhak{19}{Kosten je Erfolg zweier Varianten über Pfadwahrscheinlichkeiten vergleichen}
\abhak{20}{Pfadregeln bei unabhängigen Stufen – Fehler finden, Begründen, Anwendung, Darstellungswechsel}
\abhakgruppe{Einheit 4 · Ziehen ohne Zurücklegen und Umlegen}
\abhak{21}{Ohne Zurücklegen und Umlegen}
\abhak{22}{Ohne Zurücklegen und Umlegen – Fehler finden, Begründen, Anwendung, Darstellungswechsel}
\abhakgruppe{Einheit 5 · Mammutbäume}
\abhak{23}{Mammutbäume}
\abhak{24}{Mammutbäume – Fehler finden, Begründen, Anwendung, Darstellungswechsel}
\abhakgruppe{Einheit 6 · Situationsbäume}
\abhak{25}{Situationsbäume}
\abhak{26}{Situationsbäume – Fehler finden, Begründen, Anwendung, Darstellungswechsel}
\abhakgruppe{Einheit 7 · Term und Ereignis}
\abhak{27}{Term und Ereignis}
\abhak{28}{Term und Ereignis – Fehler finden, Begründen, Anwendung, Darstellungswechsel}
\abhakgruppe{Einheit 8 · Rückwärts}
\abhak{29}{Rückwärts}
\abhak{30}{Rückwärts – Fehler finden, Begründen, Anwendung, Darstellungswechsel}
\end{abhakseite}
